\section{Related Work}

\paragraph{Streaming speech translation.}
SST systems translate partial speech while balancing quality and latency.
Prefix-to-prefix policies make this trade-off explicit
\citep{ma2019stacl}, attention-based policies refine read/write decisions
\citep{papi2023alignatt,papi2024streamatt}, and SimulEval standardizes
streaming evaluation \citep{ma2020simuleval}. Recent speech-language models
broaden the modeling choices for direct SST
\citep{seamless2023,zhang2024streamspeech,xu2025qwen3omni}. AutoTerm-SST does
not propose a new translation model; it contributes a terminology-management
and demo layer around an existing streaming model.

\paragraph{Terminology-aware translation.}
Terminology control is commonly implemented with user-provided dictionaries,
constrained decoding, or training-time exposure to constraints
\citep{hokamp2017lexically,dinu2019training}. Terms and named entities remain a
known weakness of speech translation \citep{gaido2021moby}. Retrieval-augmented
generation offers a non-parametric alternative
\citep{lewis2020retrieval,khandelwal2021nearest}, and RASST applies glossary
retrieval before streaming LLM decoding \citep{luo2026rasst}. These approaches
normally assume that the relevant glossary is known before translation.
AutoTerm instead manages a budgeted cache of topic slices online.

\paragraph{Interactive translation systems.}
Serving engines such as vLLM~\citep{kwon2023efficient} and
SGLang~\citep{zheng2024sglang} provide continuous batching, while interactive translation tools
make model behavior inspectable \citep{dandekar2024translation}. AutoTerm-SST
combines streaming audio transport with a live terminology evidence view and
automatic glossary selection.
